\documentclass[11pt]{article}

\usepackage[margin=1in]{geometry}
\usepackage{amsmath,amssymb,amsfonts}
\usepackage{booktabs}
\usepackage{enumitem}
\usepackage{hyperref}

\hypersetup{
    colorlinks=true,
    linkcolor=blue,
    citecolor=blue,
    urlcolor=blue
}

\setlist[itemize]{leftmargin=1.5em, itemsep=0.2em, topsep=0.2em}
\setlist[enumerate]{leftmargin=1.7em, itemsep=0.25em, topsep=0.25em}

\title{Offline Metrics for Early Stopping in VLA Fine-Tuning}
\author{}
\date{}

\begin{document}
\maketitle

\section{Goal}

We want to decide when to stop offline fine-tuning of a Vision--Language--Action (VLA) policy without evaluating every checkpoint on a real robot.

For each checkpoint, we compute offline validation metrics and ask:
\begin{quote}
Which offline metric best predicts downstream execution quality?
\end{quote}

The candidate metrics are:
\begin{enumerate}
    \item Dynamic Time Warping distance: DTW;
    \item Optimal Transport distance: OT;
    \item cosine distance: COS;
    \item negative log-likelihood or a model-specific surrogate: NLL.
\end{enumerate}

All metrics are written so that \textbf{lower is better}.

\section{Reader-Friendly Notation}

We use only a few symbols.

\begin{center}
\begin{tabular}{ll}
\toprule
Symbol & Meaning \\
\midrule
$k$ & checkpoint index \\
$h_t$ & policy input at time $t$, e.g., image, robot state, and language \\
$a_t$ & robot action at time $t$ \\
$H$ & evaluation horizon, or action-chunk length \\
$f_k$ & one-step VLA at checkpoint $k$ \\
$F_k$ & chunking VLA at checkpoint $k$ \\
$W$ & world model \\
$A^*$ & expert action chunk \\
$A_k$ & generated action chunk from checkpoint $k$ \\
$R^*$ & expert state-action rollout \\
$R_k$ & generated state-action rollout from checkpoint $k$ \\
\bottomrule
\end{tabular}
\end{center}

The star $*$ means expert/reference.
The subscript $k$ means generated by checkpoint $k$.
We avoid hats and long superscripts.

For one validation window, the expert action chunk is
\begin{equation}
A^*=(a^*_0,a^*_1,\ldots,a^*_{H-1}).
\end{equation}

The generated action chunk is
\begin{equation}
A_k=(a_{k,0},a_{k,1},\ldots,a_{k,H-1}).
\end{equation}

Before computing distances, normalize each action dimension using training-set statistics:
\begin{equation}
\bar a_d = \frac{a_d-\mu_d}{\sigma_d+\epsilon}.
\end{equation}



\section{\textcolor{red}{What models to be used}}

\begin{enumerate}
    \item  \textbf{\textcolor{red}{World models:} a pretrained light-weighted model: \url{https://rlinf.readthedocs.io/en/latest/rst_source/examples/embodied/opensora.html}}


    \item \textbf{\textcolor{blue}{VLA models:} OpenVLA, Pi0.5-fast, SmoVLA}
\end{enumerate}


\section{How to Generate $A_k$ and $R_k$}

\subsection{Case 1: One-Step VLA, No World Model}

A one-step VLA predicts one action from one input:
\begin{equation}
a=f_k(h).
\end{equation}

Without a world model, feed the expert inputs into the policy one by one:
\begin{equation}
a_{k,j}=f_k(h^*_j),
\qquad j=0,1,\ldots,H-1.
\end{equation}

This gives the action chunk $A_k$.
This is a \textbf{teacher-forced} action chunk because every input comes from the expert trajectory.
It measures action quality on expert inputs, but not closed-loop error accumulation.

\subsection{Case 2: One-Step VLA, With a World Model}

With a world model, the one-step VLA is evaluated in closed loop.
The policy predicts one action, then the world model predicts the next input.

Start from the expert initial input:
\begin{equation}
h_{k,0}=h^*_0.
\end{equation}

Then repeat:
\begin{align}
a_{k,j} &= f_k(h_{k,j}), \\
h_{k,j+1} &= W(h_{k,j},a_{k,j}),
\qquad j=0,1,\ldots,H-1.
\end{align}

The generated actions form $A_k$.
The generated state-action trajectory is denoted by $R_k$.

Only the first input comes from the dataset. Future inputs come from the world model, so this setting can measure closed-loop compounding error.

\subsection{Case 3: Chunking VLA, No World Model}

A chunking VLA directly predicts the whole action chunk:
\begin{equation}
A_k=F_k(h^*_0).
\end{equation}

This is the natural action-chunk evaluation mode for chunking VLAs such as $\pi_{0.5}$-style or SmolVLA-style models.

\subsection{Case 4: Chunking VLA, With a World Model}

First generate the action chunk:
\begin{equation}
A_k=F_k(h^*_0).
\end{equation}

Then execute this action chunk in the world model:
\begin{align}
h_{k,0} &= h^*_0, \\
h_{k,j+1} &= W(h_{k,j},a_{k,j}),
\qquad j=0,1,\ldots,H-1.
\end{align}

The resulting state-action trajectory is denoted by $R_k$.

\section{What to Compare}

There are two evaluation modes.

\subsection{Mode A: Action-Chunk Evaluation}

Compare generated and expert action chunks:
\begin{equation}
A_k \quad \text{vs.} \quad A^*.
\end{equation}

Use this mode to evaluate action-level similarity.

\subsection{Mode B: State-Action Rollout Evaluation}

Use this mode only when a world model is available.

Compare generated and expert rollouts:
\begin{equation}
R_k \quad \text{vs.} \quad R^*.
\end{equation}

In implementation, do not compare raw images directly. Convert each state-action pair into a feature vector:
\begin{equation}
r_t=E(h_t,a_t),
\end{equation}
where $E$ may concatenate an image embedding, a robot-state embedding, and the normalized action.
Then apply DTW, OT, or cosine distance to the resulting feature sequences.

\section{Metric Definitions}

Let $X_k$ be the generated sequence and $X^*$ be the expert sequence.

For action-chunk evaluation:
\begin{equation}
X_k=A_k,
\qquad
X^*=A^*.
\end{equation}

For state-action rollout evaluation:
\begin{equation}
X_k=R_k,
\qquad
X^*=R^*.
\end{equation}

\subsection{DTW}

DTW is useful when the same behavior may happen at slightly different speeds:
\begin{equation}
D_{\mathrm{DTW}}(k)=\mathrm{DTW}(X_k,X^*).
\end{equation}

Use squared Euclidean distance inside DTW:
\begin{equation}
d(u,v)=\|u-v\|_2^2.
\end{equation}

\subsection{OT}

OT compares the overall distribution of points in the two sequences:
\begin{equation}
D_{\mathrm{OT}}(k)=\mathrm{SinkhornOT}(X_k,X^*).
\end{equation}

Use the same pointwise distance $d(u,v)=\|u-v\|_2^2$.

\subsection{Cosine Distance}

Define cosine similarity as
\begin{equation}
\mathrm{cos}(u,v)=\frac{u^\top v}{\|u\|_2\|v\|_2+\epsilon}.
\end{equation}

The cosine distance is
\begin{equation}
D_{\mathrm{cos}}(k)=1-\frac{1}{H}\sum_{j=0}^{H-1}\mathrm{cos}(x_{k,j},x^*_j).
\end{equation}

Cosine distance measures direction, not magnitude. It should not be used alone.

\subsection{Negative Log-Likelihood}

The correct likelihood object is
\begin{equation}
\log \pi_k(a^*_t\mid h^*_t),
\end{equation}
not $\log(a_t\mid h_t)$.

For a one-step VLA, use teacher-forced NLL:
\begin{equation}
D_{\mathrm{NLL}}(k)=-\frac{1}{H}\sum_{j=0}^{H-1}\log \pi_k(a^*_j\mid h^*_j).
\end{equation}

For a chunking VLA with tractable chunk likelihood, use
\begin{equation}
D_{\mathrm{NLL}}(k)=-\frac{1}{H}\log \pi_k(A^*\mid h^*_0).
\end{equation}

If the model does not provide a tractable likelihood, use the correct validation surrogate:
\begin{itemize}
    \item discrete action-token policy: token NLL;
    \item deterministic continuous policy: L1 or L2 validation error;
    \item diffusion or flow-matching policy: validation diffusion or flow-matching loss.
\end{itemize}

Do not use the likelihood of the policy's own generated actions as the main quality metric. A wrong but overconfident policy can assign high likelihood to its own bad actions.

\section{Aggregate over Validation Windows}

Let $m_w(k)$ be metric $m$ on validation window $w$. Average over all validation windows:
\begin{equation}
m(k)=\frac{1}{|\mathcal V|}\sum_{w\in\mathcal V}m_w(k).
\end{equation}

Use the same validation windows for every checkpoint and every metric.

\section{Early-Stopping Rule}

For metric $m$, define the best observed validation value:
\begin{equation}
m_{\mathrm{best}}=\min_k m(k).
\end{equation}

Choose the earliest checkpoint close to the best value:
\begin{equation}
k_m=\min\{k:m(k)\le m_{\mathrm{best}}+\eta\}.
\end{equation}

A simple tolerance is
\begin{equation}
\eta=0.01\,|m(0)-m_{\mathrm{best}}|.
\end{equation}

This selects the earliest checkpoint within one percent of the best observed improvement.

\section{How to Decide Which Metric Is Best}

To evaluate a metric, compare it with a trusted downstream quality score $Q(k)$.

Examples of $Q(k)$ are:
\begin{itemize}
    \item simulator success rate;
    \item small real-robot validation success rate;
    \item human progress score;
    \item validated world-model success score.
\end{itemize}

A good offline metric should rank checkpoints similarly to $Q(k)$:
\begin{equation}
\rho_m=\mathrm{SpearmanCorr}(-m(k),Q(k)).
\end{equation}

The minus sign is used because lower metric value should indicate better policy quality.

Also compute selection regret:
\begin{equation}
\mathrm{Regret}(m)=\max_k Q(k)-Q(k_m).
\end{equation}

A good metric has high $\rho_m$ and low regret.

\section{Experiment Procedure}

\begin{enumerate}
    \item Split demonstrations into training, validation, and test sets.
    \item Normalize actions using training-set statistics only.
    \item Choose the horizon $H$ and a fixed set of validation windows $\mathcal V$.
    \item Fine-tune each VLA and save checkpoints.
    \item For each checkpoint, generate $A_k$ or $R_k$.
    \item Compute DTW, OT, cosine distance, and NLL or the correct surrogate.
    \item Use each metric to select an early-stopping checkpoint.
    \item Compare the selected checkpoints using $Q(k)$, if available.
\end{enumerate}

\section{Recommended Tables}

\subsection*{Metric Curve Table}

\begin{center}
\begin{tabular}{llllr}
\toprule
Model & Checkpoint & Mode & Metric & Value \\
\midrule
OpenVLA-like & $k$ & Action-TF & DTW & \\
OpenVLA-like & $k$ & Action-TF & OT & \\
OpenVLA-like & $k$ & Action-TF & NLL & \\
Chunk VLA & $k$ & Action-Chunk & DTW & \\
Chunk VLA & $k$ & Action-Chunk & OT & \\
Chunk VLA & $k$ & Action-Chunk & NLL/surrogate & \\
Any VLA & $k$ & State-Action-WM & DTW & \\
Any VLA & $k$ & State-Action-WM & OT & \\
Any VLA & $k$ & State-Action-WM & COS & \\
\bottomrule
\end{tabular}
\end{center}

\subsection*{Metric Ranking Table}

\begin{center}
\begin{tabular}{llllrrr}
\toprule
Metric & Mode & Selected ckpt & $Q(k_m)$ & Spearman $\rho_m$ & Regret & Comment \\
\midrule
DTW & Action-TF & & & & & \\
OT & Action-TF & & & & & \\
COS & Action-TF & & & & & \\
NLL & Action-TF & & & & & \\
DTW & State-Action-WM & & & & & \\
OT & State-Action-WM & & & & & \\
COS & State-Action-WM & & & & & \\
Combined & State-Action-WM & & & & & \\
\bottomrule
\end{tabular}
\end{center}

\section{Main Claim}

The correct claim is:
\begin{quote}
Among the evaluated offline metrics, metric $m$ is the best early-stopping signal if it has the highest rank correlation with downstream quality and the lowest checkpoint-selection regret.
\end{quote}

Do not claim that an offline metric proves real-world success by itself.

The key distinction is:
\begin{equation}
\boxed{\text{teacher-forced action similarity} \neq \text{closed-loop rollout quality}.}
\end{equation}

The final principle is:
\begin{equation}
\boxed{\text{choose the earliest checkpoint selected by a validated offline proxy}.}
\end{equation}


\section{Experiment implement}

The evaluation work is split in two halves:

\textbf{Offline half (\verb|checkpoint_selection.py|)} For a fine-tuning \verb|run_dir| full of \verb|step_*| checkpoints, load each checkpoint's offline metrics (DTW, OT, cosine distance, NLL/flow-matching loss) — from \verb|metrics_history.json| if present, else fall back to wandb. For each metric, apply the early-stopping rule: pick the \textit{earliest} checkpoint within 1\% of that metric's best-ever value. This gives one candidate checkpoint per metric, e.g. \verb|{val_dtw_distance: step 5999, val_nll: step 7499, ...}|.

\textbf{Online half (the two scripts differ here)} Actually run those candidate checkpoints through a real LIBERO simulator rollout to get ground-truth success rate to evaluate does the offline metric actually predict downstream quality signal.

\begin{itemize}
    \item \textbf{\verb|openvla/.../select_and_eval_checkpoints.py|} shells out to \verb|run_libero_eval.py| (subprocess, once per selected checkpoint), reading back the summary JSON it writes.
    \item \textbf{\verb|lerobot/.../select_and_eval_checkpoints.py|} shells out to lerobot's own \verb|lerobot-eval| CLI, which already knows how to load a LoRA/PEFT \verb|pi05| checkpoint on \verb|LiberoEnv| — no custom eval loop needed there.
\end{itemize}
\textbf{Output} cache each checkpoint's result, then print/save the tex's Metric Ranking Table: for each metric, its Spearman correlation with the real success rate and its regret value, which calculate how much success rate the metric's selected checkpoint gave up compare to the actual best one.
\end{document}
